
\documentclass{article}
\usepackage{colortbl}
\usepackage{makecell}
\usepackage{multirow}
\usepackage{supertabular}

\begin{document}

\newcounter{utterance}

\centering \large Interaction Transcript for game `cladder', experiment `full\_v1.5\_default', episode 4315 with LLP: Large Language Problems\_\_qwen3.5{-}9b\_\_llp{-}final.
\vspace{24pt}

{ \footnotesize  \setcounter{utterance}{1}
\setlength{\tabcolsep}{0pt}
\begin{supertabular}{c@{$\;$}|p{.15\linewidth}@{}p{.15\linewidth}p{.15\linewidth}p{.15\linewidth}p{.15\linewidth}p{.15\linewidth}}
    \# & \multicolumn{2}{c}{Player} && \multicolumn{2}{c}{Game Master} \\
    \hline

    \theutterance \stepcounter{utterance}  
    & & & \multicolumn{4}{p{0.6\linewidth}}{
        \cellcolor[rgb]{0.9,0.9,0.9}{
            \makecell[{{p{\linewidth}}}]{
                \texttt{\tiny{[P1$\langle$GM]}}
                \texttt{You are an expert in causal inference. The following question is not a typical commonsense query, but rather a meticulously designed question created by a professor specializing in causal inference, intended to assess the students' mastery of the course content.} \\
\texttt{Imagine a self{-}contained, hypothetical world with only the following conditions, and without any unmentioned factors or causal relationships: Maternal smoking status has a direct effect on infant's birth weight and freckles. Health condition has a direct effect on infant's birth weight and freckles. Infant's birth weight has a direct effect on freckles. Health condition is unobserved. For infants with nonsmoking mothers and low infant birth weight, the probability of freckles is 46\%. For infants with nonsmoking mothers and normal infant birth weight, the probability of freckles is 21\%. For infants with smoking mothers and low infant birth weight, the probability of freckles is 74\%. For infants with smoking mothers and normal infant birth weight, the probability of freckles is 47\%. For infants with nonsmoking mothers and with poor health, the probability of normal infant birth weight is 51\%. For infants with nonsmoking mothers and with good health, the probability of normal infant birth weight is 72\%. For infants with smoking mothers and with poor health, the probability of normal infant birth weight is 3\%. For infants with smoking mothers and with good health, the probability of normal infant birth weight is 35\%. The overall probability of good health is 90\%. If we disregard the mediation effect through infant's birth weight, would maternal smoking status negatively affect freckles?} \\
\texttt{Start your answer with 'Yes' or 'No', followed by additional reasoning or evidence to support your explanation.} \\
            }
        }
    }
     \\ \\

    \theutterance \stepcounter{utterance}  
    & \multicolumn{4}{p{0.6\linewidth}}{
        \cellcolor[rgb]{0.561,0.737,0.561}{
            \makecell[{{p{\linewidth}}}]{
                \texttt{\tiny{[P1$\rangle$GM]}}
                \texttt{No.} \\
\\ 
\texttt{To determine whether maternal smoking status negatively affects freckles when disregarding the mediation effect through infant's birth weight, we need to compare the probability of freckles for smoking versus nonsmoking mothers while holding health condition constant (since health is unobserved but affects both birth weight and freckles).} \\
\\ 
\texttt{Let's calculate the conditional probabilities of freckles given maternal smoking status and health condition:} \\
\\ 
\texttt{For nonsmoking mothers:} \\
\texttt{{-} Given poor health: P(freckles | nonsmoking, poor health) = 46\% (since low birth weight is the only path, and we're disregarding birth weight mediation)} \\
\texttt{{-} Given good health: P(freckles | nonsmoking, good health) = 21\% (since normal birth weight is the only path, and we're disregarding birth weight mediation)} \\
\\ 
\texttt{For smoking mothers:} \\
\texttt{{-} Given poor health: P(freckles | smoking, poor health) = 74\%} \\
\texttt{{-} Given good health: P(freckles | smoking, good health) = 47\%} \\
\\ 
\texttt{Now comparing at each health level:} \\
\texttt{{-} Poor health: 74\% (smoking) vs 46\% (nonsmoking) {-} smoking increases freckles} \\
\texttt{{-} Good health: 47\% (smoking) vs 21\% (nonsmoking) {-} smoking increases freckles} \\
\\ 
\texttt{In both health conditions, smoking mothers have higher probability of freckles than nonsmoking mothers. Therefore, maternal smoking status does NOT negatively affect freckles; it positively affects freckles when disregarding the mediation through birth weight.} \\
            }
        }
    }
    & & \\ \\

    \theutterance \stepcounter{utterance}  
    & & & \multicolumn{2}{p{0.3\linewidth}}{
        \cellcolor[rgb]{0.9,0.9,0.9}{
            \makecell[{{p{\linewidth}}}]{
                \texttt{\tiny{[GM$|$GM]}}
                \texttt{No} \\
            }
        }
    }
    & & \\ \\

    \theutterance \stepcounter{utterance}  
    & & & \multicolumn{2}{p{0.3\linewidth}}{
        \cellcolor[rgb]{0.9,0.9,0.9}{
            \makecell[{{p{\linewidth}}}]{
                \texttt{\tiny{[GM$|$GM]}}
                \texttt{no} \\
            }
        }
    }
    & & \\ \\

    \theutterance \stepcounter{utterance}  
    & & & \multicolumn{2}{p{0.3\linewidth}}{
        \cellcolor[rgb]{0.9,0.9,0.9}{
            \makecell[{{p{\linewidth}}}]{
                \texttt{\tiny{[GM$|$GM]}}
                \texttt{game\_result = WIN} \\
            }
        }
    }
    & & \\ \\

\end{supertabular}
}

\end{document}
